\documentclass[11pt, a4paper]{article}

% --- PAQUETES PRINCIPALES ---
\usepackage[utf8]{inputenc}
\usepackage[spanish, es-tabla]{babel}
\usepackage[margin=2.5cm]{geometry}
\usepackage{amsmath, amssymb, mathtools}
\usepackage{siunitx}
\usepackage{graphicx}
\usepackage{booktabs}
\usepackage{tabularx}
\usepackage[most]{tcolorbox}
\usepackage{listings}
\usepackage{xcolor}

% --- PAQUETES PARA GRÁFICAS Y CIRCUITOS ---
\usepackage{pgfplots}
\usepackage[american]{circuitikz}
\pgfplotsset{compat=1.18}

% --- CONFIGURACIÓN DE CÓDIGO (PYTHON) ---
\definecolor{codegreen}{rgb}{0,0.6,0}
\definecolor{codegray}{rgb}{0.5,0.5,0.5}
\definecolor{codepurple}{rgb}{0.58,0,0.82}
\definecolor{backcolour}{rgb}{0.96,0.96,0.96}
\lstset{
    backgroundcolor=\color{backcolour},   
    commentstyle=\color{codegreen},
    keywordstyle=\color{magenta},
    numberstyle=\tiny\color{codegray},
    stringstyle=\color{codepurple},
    basicstyle=\ttfamily\footnotesize,
    breaklines=true,                 
    numbers=left,                    
    numbersep=5pt,                  
    language=Python,
    literate={á}{{\'a}}1 {é}{{\'e}}1 {í}{{\'i}}1 {ó}{{\'o}}1 {ú}{{\'u}}1
             {Á}{{\'A}}1 {É}{{\'E}}1 {Í}{{\'I}}1 {Ó}{{\'O}}1 {Ú}{{\'U}}1
             {ñ}{{\~n}}1 {Ñ}{{\~N}}1
}

% --- CONFIGURACIÓN DE SIUNITX ---
\sisetup{output-decimal-marker = {,}, inter-unit-product = \ensuremath{{}\cdot{}}}

% --- CABECERAS ---
\usepackage{fancyhdr}
\pagestyle{fancy}
\fancyhf{}
\lhead{\textbf{Circuitos Electrónicos III (IMT-221)}}
\rhead{Sesión 5: Rectificación y Filtrado Capacitivo}
\cfoot{\thepage}
\renewcommand{\headrulewidth}{0.4pt}

% --- ENTORNOS PERSONALIZADOS ---
\newtcolorbox{aviso}[1][]{colback=red!5!white, colframe=red!75!black, fonttitle=\bfseries, title=#1}

\begin{document}

\begin{center}
    {\LARGE \textbf{Apuntes de Cátedra: Sesión 5}}\\[0.3cm]
    {\large \textbf{Rectificación Monofásica (Puente Completo) y Filtrado Capacitivo}}\\
    \textit{Docente: M.Sc. Ing. Bernardo Quiroga Turdera}
    \vspace{0.3cm}
    \hrule
\end{center}

\vspace{0.3cm}

% ==========================================
% SECCIÓN 1: FUNDAMENTACIÓN FÍSICA Y TOPOLOGÍAS
% ==========================================
\section{Fundamentación Física y Topologías}
Todo actuador mecatrónico o equipo electromédico requiere un bus de corriente continua estable a partir de la red de CA. La rectificación sin filtro entrega una tensión pulsante inutilizable para alimentar etapas de control o instrumentación (como en el Hito 2 del PFI). 

\subsection{Comparativa de Topologías de Rectificación}
\begin{center}
\renewcommand{\arraystretch}{1.5}
\begin{tabularx}{\textwidth}{@{} >{\bfseries}l X X X @{}}
\toprule
Parámetro de Diseño & Media Onda & Onda Completa (Tap) & Puente de Graetz \\
\midrule
Diodos requeridos & 1 & 2 & 4 \\
Tensión pico de salida ($V_{p(\text{out})}$) & $V_p - V_\gamma$ & $V_p - V_\gamma$ & $V_p - 2V_\gamma$ \\
Frecuencia de rizado ($f_r$) & $f_{in}$ (\qty{50}{\hertz}) & $2f_{in}$ (\qty{100}{\hertz}) & $2f_{in}$ (\qty{100}{\hertz}) \\
Tensión Inversa de Pico (PIV) & $2V_p$ & $2V_p$ & $V_p$ \\
Eficiencia de conversión & $40.6\%$ & $81.2\%$ & $81.2\%$ \\
\bottomrule
\end{tabularx}
\end{center}

\subsection{Mecánica del Filtrado Capacitivo}
Durante el primer cuarto de ciclo, el diodo conduce y carga el capacitor hasta el valor pico de salida $V_{p(\text{out})}$. Cuando la tensión de la fuente senoidal cae por debajo de la tensión almacenada en el capacitor, los diodos se polarizan en inversa (se apagan). El capacitor entrega energía a la carga $R_L$ mediante una descarga exponencial:
\begin{equation}
    v_o(t) = V_{p(\text{out})} \cdot e^{-\frac{t}{R_L C}}
\end{equation}
Si la constante de tiempo $\tau = R_L C \gg T$, la descarga se aproxima a una recta.

\begin{figure}[h!]
    \centering
    \begin{tikzpicture}[scale=0.9]
        \begin{axis}[
            width=12cm, height=6cm,
            axis lines=middle,
            xmin=0, xmax=10, ymin=0, ymax=1.2,
            xtick={1.57, 4.71, 7.85}, xticklabels={$\frac{T}{4}$,$\frac{3T}{4}$,$\frac{5T}{4}$},
            ytick={1, 0.7}, yticklabels={$V_{p(\mathrm{out})}$,$V_{p}-V_{r}$},
            xlabel={$t$}, ylabel={$v_o(t)$},
            grid=both, grid style={dashed, gray!30},
            legend style={at={(0.5,1.1)}, anchor=south, legend columns=-1}
        ]
        % Onda rectificada completa (Valor absoluto del seno)
        \addplot[domain=0:10, samples=200, gray, dashed, thick] {abs(sin(deg(x)))};
        \addlegendentry{CA Rectificada (Sin filtro)}
        
        % Carga y Descarga Capacitiva
        \addplot[domain=0:1.57, samples=50, blue, ultra thick] {sin(deg(x))};
        \addplot[domain=1.57:3.7, samples=50, blue, ultra thick] {exp(-0.15*(x-1.57))};
        \addplot[domain=3.7:4.71, samples=50, blue, ultra thick] {abs(sin(deg(x)))};
        \addplot[domain=4.71:6.84, samples=50, blue, ultra thick] {exp(-0.15*(x-4.71))};
        \addplot[domain=6.84:7.85, samples=50, blue, ultra thick] {abs(sin(deg(x)))};
        \addlegendentry{Voltaje con Filtro $C$}
        
        % Indicador de Rizado
        \draw[<->, red, thick] (axis cs: 6.84, 0.72) -- (axis cs: 6.84, 1) node[midway, right] {$V_r$};
        \end{axis}
    \end{tikzpicture}
    \caption{Respuesta de voltaje en un rectificador de onda completa con filtro capacitivo.}
\end{figure}

% ==========================================
% SECCIÓN 2: DEDUCCIÓN ANALÍTICA
% ==========================================
\section{Deducción Analítica del Rizado y Corriente Pico}
\textbf{1. Tensión de Rizado ($V_r$) y Tensión en Continua ($V_{DC}$):}
Para un rectificador de onda completa, la tensión de rizado pico a pico se calcula como:
\begin{equation}
    V_r \approx \frac{I_L}{2 f_{in} C} = \frac{V_{p(\text{out})}}{2 f_{in} R_L C}
\end{equation}
La tensión media de continua efectiva en la carga es:
\begin{equation}
    V_{DC} = V_{p(\text{out})} - \frac{V_r}{2}
\end{equation}

\textbf{2. Ángulo de Conducción ($\theta_c$) y Corriente Pico ($I_{D(\max)}$):}\\
El diodo empieza a conducir cuando la senoidal iguala la tensión del capacitor en el punto de valle ($V_p - V_r$):
\begin{equation}
    V_p \cos(\theta_c) = V_p - V_r \implies \cos(\theta_c) = 1 - \frac{V_r}{V_p}
\end{equation}
Aproximando $\cos(\theta_c) \approx 1 - \frac{\theta_c^2}{2}$ por serie de Maclaurin (para $\theta_c \ll 1$):
\begin{equation}
    1 - \frac{\theta_c^2}{2} = 1 - \frac{V_r}{V_p} \implies \theta_c = \sqrt{\frac{2 V_r}{V_p}}
\end{equation}
El tiempo de conducción es sumamente breve: $\Delta t = \frac{\theta_c}{\omega} = \frac{1}{2\pi f} \sqrt{\frac{2 V_r}{V_p}}$. 

Por el principio de \textbf{Conservación de Carga}, la carga extraída en un semiciclo ($\Delta Q_{out} = I_L \cdot \frac{T}{2}$) debe reponerse íntegramente durante el breve pulso de conducción $\Delta t$. Esto genera una elevadísima corriente pico en el diodo:
\begin{equation}
    I_{D(\max)} \approx I_L \left( 1 + \pi \sqrt{\frac{2 V_{p(\text{out})}}{V_r}} \right)
\end{equation}
\textit{Costo oculto:} A mayor capacitancia $C$, menor es el rizado $V_r$, pero mayor es el impacto de la corriente $I_{D(\max)}$ que puede destruir la unión P-N del diodo.

\newpage
% ==========================================
% SECCIÓN 3: PROBLEMA DE DISEÑO (PFI)
% ==========================================
\section{Problema de Diseño de Ingeniería (PFI)}
\textbf{Enunciado:} Diseñar la etapa rectificadora para la fuente de instrumentación del PFI. Se dispone de un transformador de $12\text{ V}_{\text{rms}}$ a $\qty{50}{\hertz}$ y un puente rectificador integrado (Graetz) con diodos de silicio ($V_\gamma = \qty{0.7}{\volt}$). La carga consume $I_L = \qty{500}{\milli\ampere}$.

\vspace{0.3cm}
\begin{center}
    \begin{circuitikz}[american, scale=0.9, transform shape]
        % Transformador / Fuente AC
        \draw (-2,0) to[sV, l=\qty{12}{\volt_{\text{rms}}}] (-2,-4);
        \draw (-2,0) |- (0,0);
        \draw (-2,-4) |- (4,-4) -- (4,0);
        
        % Puente de Graetz (Diodos)
        \draw (0,0) to[D, l=$D_2$, *-*] (2,2);
        \draw (4,0) to[D, l_=$D_1$, *-*] (2,2);
        \draw (2,-2) to[D, l=$D_3$, *-*] (0,0);
        \draw (2,-2) to[D, l_=$D_4$, *-*] (4,0);
        
        % Bus DC, Capacitor y Carga
        \draw (2,2) -- (6,2) to[C, l=$C$] (6,-2) -- (2,-2);
        \draw (6,2) -- (8,2) to[R, l=$R_L$ (\qty{500}{\milli\ampere})] (8,-2) -- (6,-2);
        \node[ground] at (7,-2) {};
    \end{circuitikz}
\end{center}

\textbf{Requerimientos:} Tensión de rizado máximo $V_r \le \qty{0.5}{\volt}$. Determinar el capacitor $C$, el voltaje continuo $V_{DC}$, la corriente pico repetitiva $I_{D(\max)}$ y seleccionar entre el 1N4148 ($I_{\max} = \qty{200}{\milli\ampere}$) o el 1N4007 ($I_{\max} = \qty{1}{\ampere}$).

\vspace{0.3cm}
\noindent \textbf{Paso 1: Tensión Pico del Secundario y Salida}\\
$V_p = 12\text{ V}_{\text{rms}} \cdot \sqrt{2} = \mathbf{\qty{16.97}{\volt}}$. Como conducen dos diodos en serie por semiciclo:
\begin{equation*}
    V_{p(\text{out})} = V_p - 2V_\gamma = \qty{16.97}{\volt} - 2(\qty{0.7}{\volt}) = \mathbf{\qty{15.57}{\volt}}
\end{equation*}

\noindent \textbf{Paso 2: Dimensionamiento del Capacitor de Filtro}\\
Para puente completo con $f_{in} = \qty{50}{\hertz} \implies f_r = \qty{100}{\hertz}$:
\begin{equation*}
    C = \frac{I_L}{f_r \cdot V_r} = \frac{\qty{0.5}{\ampere}}{\qty{100}{\hertz} \cdot \qty{0.5}{\volt}} = \qty{0.010}{\farad} = \mathbf{10\,000\,\mu\text{F}}
\end{equation*}
\textit{Selección Comercial:} Un capacitor electrolítico de $10\,000\,\mu\text{F} / 25\text{V}$.

\noindent \textbf{Paso 3: Tensión Continua Promedio ($V_{DC}$)}\\
\begin{equation*}
    V_{DC} = V_{p(\text{out})} - \frac{V_r}{2} = 15.57\text{ V} - \frac{0.5\text{ V}}{2} = \mathbf{\qty{15.32}{\volt}}
\end{equation*}

\noindent \textbf{Paso 4: Corriente Pico Repetitiva en los Diodos ($I_{D(\max)}$)}\\
Calculamos el ángulo de conducción:
\begin{equation*}
    \theta_c = \sqrt{\frac{2 \cdot 0.5\text{ V}}{15.57\text{ V}}} = \sqrt{0.0642} = \qty{0.2534}{\radian} \approx 14.52^\circ
\end{equation*}
Calculamos la corriente pico que atraviesa los semiconductores:
\begin{equation*}
    I_{D(\max)} = I_L \left( 1 + \pi \sqrt{\frac{2 \cdot 15.57}{0.5}} \right) = 0.5 \left( 1 + \pi \sqrt{62.28} \right) = \mathbf{\qty{12.89}{\ampere}}
\end{equation*}
\textit{Nota Técnica:} Aunque la carga consume solo $\qty{500}{\milli\ampere}$, los diodos sufren pulsos de casi $\qty{13}{\ampere}$ repetidamente cada $\qty{10}{\milli\second}$.

\noindent \textbf{Paso 5: Selección de Componente y PIV}\\
La Tensión Inversa de Pico soportada por los diodos apagados en el puente es:
\begin{equation*}
    PIV = V_p - V_\gamma = 16.97\text{ V} - 0.7\text{ V} = \mathbf{\qty{16.27}{\volt}}
\end{equation*}
El diodo 1N4148 ($I_{\max} = \qty{200}{\milli\ampere}$) se destruiría inmediatamente debido al pulso de $\qty{13}{\ampere}$. La \textbf{selección obligatoria es el 1N4007}, que soporta \qty{1}{\ampere} promedio, pero tiene una capacidad de pulso transitorio no repetitivo ($I_{FSM}$) de \qty{30}{\ampere} y un PIV de \qty{1000}{\volt}.

\newpage
% ==========================================
% SECCIÓN 4: PYTHON Y AUDITORÍA
% ==========================================
\section{Taller Computacional: Descarga Exponencial y Diodos}
Este script compara la aproximación lineal de rizado vs. la solución exacta con descarga exponencial.

\begin{lstlisting}
# ==============================================================================
# UNIVERSIDAD CATOLICA BOLIVIANA "SAN PABLO" - SEDE TARIJA
# Sesion 5: Modelado No Lineal de Rectificador Puente y Filtro Capacitivo
# Docente: M.Sc. Ing. Bernardo Quiroga Turdera
# ==============================================================================
import numpy as np
import matplotlib.pyplot as plt

# 1. PARAMETROS DEL SISTEMA
f_in = 50.0           # Frecuencia de red (Hz)
V_rms = 12.0          # Voltaje eficaz del secundario (V)
V_gamma = 0.7         # Caida directa por diodo (V)
I_load_nominal = 0.5  # Corriente DC de carga (A)
V_ripple_spec = 0.5   # Rizado pico a pico deseado (V)

V_peak = V_rms * np.sqrt(2)
V_peak_out = V_peak - 2 * V_gamma
R_load = V_peak_out / I_load_nominal

# Dimensionamiento Analitico
f_ripple = 2 * f_in
C_filtro = I_load_nominal / (f_ripple * V_ripple_spec)

print(f"Tension Pico Salida Rectificada : {V_peak_out:.2f} V")
print(f"Capacitancia Requerida (C)      : {C_filtro*1e6:.1f} uF")

# 2. SIMULACION TEMPORAL DINAMICA
num_ciclos = 3
t = np.linspace(0, num_ciclos / f_in, 5000)
dt = t[1] - t[0]

v_in = V_peak * np.sin(2 * np.pi * f_in * t)
v_rect = np.maximum(np.abs(v_in) - 2 * V_gamma, 0.0)

v_out = np.zeros_like(t)
i_diode = np.zeros_like(t)
v_cap = 0.0

for i in range(len(t)):
    if v_rect[i] > v_cap:
        # Conduccion: Carga del capacitor
        v_cap = v_rect[i]
        if i > 0:
            i_cap = C_filtro * (v_rect[i] - v_rect[i-1]) / dt
        else:
            i_cap = 0.0
        i_diode[i] = np.maximum(i_cap + v_cap / R_load, 0.0)
    else:
        # Bloqueo: Descarga exponencial a traves de RL
        v_cap = v_cap * np.exp(-dt / (R_load * C_filtro))
        i_diode[i] = 0.0
    v_out[i] = v_cap

# 3. GRAFICACION
plt.figure(figsize=(10, 7))

plt.subplot(2, 1, 1)
plt.plot(t * 1e3, v_rect, 'gray', linestyle='--', label='Sin Filtro')
plt.plot(t * 1e3, v_out, 'b-', linewidth=2, label=f'Filtrada (C={C_filtro*1e6:.0f} uF)')
plt.title('Respuesta Temporal: Rectificador con Filtro')
plt.ylabel('Tension (V)')
plt.grid(True, linestyle='--', alpha=0.7)
plt.legend()

plt.subplot(2, 1, 2)
plt.plot(t * 1e3, i_diode, 'r-', linewidth=1.5, label='Pulsos de Corriente ID')
plt.title('Pulsos de Corriente en Diodos')
plt.xlabel('Tiempo (ms)')
plt.ylabel('Corriente (A)')
plt.grid(True, linestyle='--', alpha=0.7)
plt.legend()

plt.tight_layout()
plt.show()
\end{lstlisting}

\begin{aviso}[Protocolo de Auditoría Dinámica (IA)]
\textbf{Consigna para el Estudiante:} En el script, cambiar el capacitor de $10\,000\,\mu\text{F}$ a $1\,000\,\mu\text{F}$ y luego a $100\,\mu\text{F}$. \\
\textbf{Pregunta:} ¿Por qué al reducir $C$ a $100\,\mu\text{F}$ la corriente pico $I_{D(\max)}$ cae de \qty{12.8}{\ampere} a menos de \qty{2}{\ampere}? \\
\textbf{Respuesta Esperada:} Porque el rizado aumenta exponencialmente. Al haber un valle más profundo, el ángulo de conducción $\theta_c$ es mucho más ancho. Al tener el diodo más tiempo ($\Delta t$) para reponer la carga, no necesita concentrar toda la energía en un pulso destructivo estrecho.
\end{aviso}

\end{document}